\documentclass{amsart}
\usepackage{amsmath,amssymb,amsthm,hyperref}

\newtheorem{theorem}{Theorem}
\newtheorem{proposition}{Proposition}
\newtheorem{lemma}{Lemma}
\newtheorem{corollary}{Corollary}
\theoremstyle{definition}
\newtheorem{definition}{Definition}
\newtheorem{remark}{Remark}

\DeclareMathOperator{\E}{\mathbb{E}}
\newcommand{\R}{\mathbb{R}}
\newcommand{\N}{\mathbb{N}}
\newcommand{\one}{\mathbf{1}}

\begin{document}

\title{Resolving the zero-prior obstruction to confirming universal hypotheses}
\maketitle

\begin{abstract}
We give a complete and rigorous resolution of the zero-prior obstruction. We
prove that the obstruction is \emph{exactly} an artifact of restricting to
priors that are absolutely continuous with respect to Lebesgue measure on the
parameter space, and we characterize the class of priors for which the
universal hypothesis $H=\{\theta=1\}$ is confirmable. The admissible class is
the class of \emph{spike-and-slab} (mixed discrete--continuous) priors that put
a strictly positive atom on $\theta=1$. For every prior in this class we show
that (i) $P(H)>0$, (ii) the posterior probability of $H$ after $n$ consecutive
successes is strictly positive and increases to $1$ as $n\to\infty$, and (iii)
the framework is an ordinary Bayesian model (hence coherent under conditioning),
is invariant under reparametrizations of $\theta$, and reduces to the
Bayes--Laplace predictive rule in the non-universal regime. We then prove a
converse: an exchangeable model confirms $H$ in the sense of (ii) if and only if
its de~Finetti mixing measure has an atom at $\theta=1$. Finally we exhibit the
algorithmic-probability (Solomonoff) realization of the same construction and
show it satisfies (i)--(iii), explaining why a Solomonoff-style universal prior
escapes the obstruction automatically.
\end{abstract}

\section{Setup and the obstruction restated}\label{sec:setup}

\subsection{The model}
Fix the parameter space $\Theta=[0,1]$. For $\theta\in\Theta$ let $P_\theta$ be
the law on $\{0,1\}^\N$ under which $X_1,X_2,\dots$ are i.i.d.\ Bernoulli$(\theta)$;
thus for every $n$ and every $x_1,\dots,x_n\in\{0,1\}$,
\begin{equation}\label{eq:likelihood}
  P_\theta(X_1=x_1,\dots,X_n=x_n)=\theta^{\,s_n}(1-\theta)^{\,n-s_n},
  \qquad s_n:=\sum_{i=1}^n x_i .
\end{equation}
A prior is a Borel probability measure $\mu$ on $\Theta$. The Bayesian model it
generates is the measure
\begin{equation}\label{eq:joint}
  P(\,\cdot\,)=\int_{\Theta}P_\theta(\,\cdot\,)\,d\mu(\theta)
\end{equation}
on $\Theta\times\{0,1\}^\N$ (with $\theta$ the first coordinate). This is the
standard exchangeable Bayesian model; by de~Finetti's theorem every
exchangeable law on $\{0,1\}^\N$ arises this way for a unique $\mu$
(Theorem~\ref{thm:converse} below uses this).

The universal hypothesis is the event
\begin{equation}\label{eq:H}
  H=\{\theta=1\}\subseteq\Theta .
\end{equation}
We write $A_n=\{X_1=\dots=X_n=1\}$ for the event of $n$ consecutive successes,
and $A_\infty=\bigcap_{n\ge1}A_n=\{X_i=1\ \forall i\}$.

\subsection{The obstruction}
We restate the problem's premise precisely, to fix exactly what must be
circumvented.

\begin{proposition}[The zero-prior obstruction]\label{prop:obstruction}
If $\mu$ is absolutely continuous with respect to Lebesgue measure on $[0,1]$,
then $P(H)=0$ and $P(H\mid A_n)=0$ for every finite $n$.
\end{proposition}

\begin{proof}
If $\mu\ll\mathrm{Leb}$ then $\mu(\{1\})=0$ because $\{1\}$ is Lebesgue-null;
hence $P(H)=\mu(H)=\mu(\{1\})=0$. For the posterior, by~\eqref{eq:joint}
and~\eqref{eq:likelihood},
\[
  P(A_n)=\int_0^1\theta^{\,n}\,d\mu(\theta),\qquad
  P(H\cap A_n)=\int_{\{1\}}\theta^{\,n}\,d\mu(\theta)=1^n\cdot\mu(\{1\})=0 .
\]
Since $\mu\ll\mathrm{Leb}$ is not the point mass at $0$, $\mu$ charges some
subinterval of $(0,1]$, so $P(A_n)=\int_0^1\theta^n\,d\mu(\theta)>0$ and
conditioning is well defined; thus
$P(H\mid A_n)=P(H\cap A_n)/P(A_n)=0$.
\end{proof}

The proof isolates the sole mechanism of the obstruction: $\mu(\{1\})=0$. The
likelihood ratio in favour of $H$ over any $\theta<1$ after $n$ successes is
$1/\theta^n\to\infty$, so the data are maximally favourable to $H$; the only
reason the posterior of $H$ cannot rise is that the \emph{prior} assigns it
measure zero. In this model all conditioning is on the
positive-probability events $A_n$, and an event of prior probability $0$ that is
contained in every conditioning event retains conditional probability $0$ under
such updating; so here no finite run of successes can lift $H$.
(This is special to conditioning on positive-probability events: in general,
disintegration of a measure along a $\sigma$-field can assign positive
conditional probability to prior-null events, so the slogan ``updating never
lifts a null event'' is \emph{not} a universal truth, only a feature of the
present setup.) The cure is therefore forced: give $H$ positive prior mass. The content
of the theorems below is that doing so can be done while preserving every
desirable feature of the continuous model.

\section{The admissible class: spike-and-slab priors}\label{sec:class}

\begin{definition}[Spike-and-slab prior]\label{def:spikeslab}
A \emph{spike-and-slab prior} is a measure of the form
\begin{equation}\label{eq:prior}
  \mu = p\,\delta_{1} + (1-p)\,\nu,
\end{equation}
where $p\in(0,1)$, $\delta_1$ is the unit point mass at $\theta=1$, and $\nu$ is
a Borel probability measure on $[0,1]$ with $\nu(\{1\})=0$ (the ``slab'').
We call $p$ the \emph{spike weight} and $\nu$ the \emph{slab law}. We say the
slab is \emph{non-degenerate} if $\nu([1-\varepsilon,1))>0$ for every
$\varepsilon>0$, i.e.\ $1\in\operatorname{supp}\nu$; the canonical choice is
$\nu=\mathrm{Beta}(a,b)$ with $a,b>0$, and in particular the
Bayes--Laplace slab $\nu=\mathrm{Unif}[0,1]$.
\end{definition}

Throughout, $P$ denotes the model~\eqref{eq:joint} built from~\eqref{eq:prior}.
This is an ordinary prior (a probability measure on $[0,1]$), so the entire
machinery of Bayesian conditioning applies verbatim; in particular requirement
(iii)'s coherence is automatic. We record the elementary integrals once.

\begin{lemma}[Marginal likelihood of a success run]\label{lem:marginal}
For every $n\ge0$,
\begin{equation}\label{eq:Pan}
  P(A_n)=p+(1-p)\,m_n,\qquad m_n:=\int_{[0,1)}\theta^{\,n}\,d\nu(\theta),
\end{equation}
with the convention $A_0$ the whole space and $m_0=1$. Moreover $m_n$ is
non-increasing in $n$ and $\lim_{n\to\infty}m_n=0$.
\end{lemma}

\begin{proof}
By~\eqref{eq:joint}, \eqref{eq:likelihood} and~\eqref{eq:prior},
\[
  P(A_n)=\int_{[0,1]}\theta^n\,d\mu(\theta)
  =p\cdot 1^n+(1-p)\int_{[0,1]}\theta^n\,d\nu(\theta).
\]
Since $\nu(\{1\})=0$, the last integral equals $\int_{[0,1)}\theta^n\,d\nu=m_n$,
giving~\eqref{eq:Pan}. For $\theta\in[0,1)$ the sequence $\theta^n$ is
non-increasing in $n$, so $m_n$ is non-increasing by monotonicity of the
integral. Finally $\theta^n\to 0$ pointwise on $[0,1)$ and $0\le\theta^n\le1$,
so by dominated convergence (dominating function $\one$, which is
$\nu$-integrable) $m_n\to\int_{[0,1)}0\,d\nu=0$.
\end{proof}

\section{Confirmation: properties (i) and (ii)}\label{sec:confirm}

\begin{theorem}[Positive prior and convergent posterior]\label{thm:main}
Let $\mu=p\,\delta_1+(1-p)\,\nu$ be any spike-and-slab prior
(Definition~\ref{def:spikeslab}). Then:
\begin{enumerate}
\item[\textup{(i)}] $P(H)=p>0$.
\item[\textup{(ii)}] For every $n\ge0$ the posterior of $H$ given $n$ consecutive
successes is
\begin{equation}\label{eq:posterior}
  P(H\mid A_n)=\frac{p}{p+(1-p)\,m_n}\in(0,1],
\end{equation}
this quantity is non-decreasing in $n$, and
\[
  \lim_{n\to\infty}P(H\mid A_n)=1.
\]
\item[\textup{(iii)}] Conditioned on the full success sequence,
$P(H\mid A_\infty)=1$; i.e.\ the model gives the correct answer in the limit.
\end{enumerate}
\end{theorem}

\begin{proof}
(i) $P(H)=\mu(\{1\})=p\,\delta_1(\{1\})+(1-p)\,\nu(\{1\})=p+0=p>0$.

(ii) The event $A_n$ has $P(A_n)=p+(1-p)m_n\ge p>0$ by
Lemma~\ref{lem:marginal}, so the conditional probability is defined. The
contribution of $H$ to $A_n$ is the spike, whose likelihood is $1^n=1$:
\[
  P(H\cap A_n)=\int_{\{1\}}\theta^n\,d\mu(\theta)=p\cdot 1^n=p .
\]
Therefore
\[
  P(H\mid A_n)=\frac{P(H\cap A_n)}{P(A_n)}=\frac{p}{p+(1-p)\,m_n},
\]
which is~\eqref{eq:posterior}. Because $p>0$ and $(1-p)m_n\ge0$ the value lies
in $(0,1]$. The map $t\mapsto p/(p+(1-p)t)$ is strictly decreasing in
$t\ge0$, and $m_n$ is non-increasing in $n$ (Lemma~\ref{lem:marginal}); hence
$P(H\mid A_n)$ is non-decreasing in $n$. Finally, since $m_n\to0$,
\[
  \lim_{n\to\infty}P(H\mid A_n)=\frac{p}{p+(1-p)\cdot 0}=\frac{p}{p}=1 .
\]

(iii) The events $A_n$ decrease to $A_\infty$, and by~\eqref{eq:joint}
\[
  P(H\cap A_\infty)=\int_{\{1\}}P_\theta(A_\infty)\,d\mu(\theta)=p\cdot 1=p,
\]
because $P_1(A_\infty)=1$. Also
$P(A_\infty)=\int_{[0,1]}P_\theta(A_\infty)\,d\mu
=p\cdot1+(1-p)\int_{[0,1)}\big(\lim_n\theta^n\big)\,d\nu$, and
$\lim_n\theta^n=0$ for $\theta\in[0,1)$ while the point $\theta=1$ is
excluded by $\nu(\{1\})=0$; by dominated convergence the slab integral is $0$, so
$P(A_\infty)=p$. Hence $P(H\mid A_\infty)=p/p=1$. (Consistently, by downward
continuity of measure $P(H\mid A_\infty)=\lim_n P(H\mid A_n)=1$.)
\end{proof}

This settles (i) and (ii) for the entire admissible class, with no restriction
on $\nu$ beyond $\nu(\{1\})=0$: the slab may be uniform, any Beta law, or any
other continuous prior on $[0,1)$. In particular the spike weight $p$ may be
taken arbitrarily small, so the resolution does not require committing strong
prior belief in $H$.

\section{Coherence, reparametrization invariance, and predictive behaviour: (iii)}
\label{sec:coherence}

\subsection{Coherence under updating}
The model $P$ of~\eqref{eq:joint} is a genuine probability measure on
$\Theta\times\{0,1\}^\N$ and $H$ a genuine event with $P(H)=p>0$. All the
posteriors above are ordinary conditional probabilities of positive-probability
events, so Bayes coherence (the chain rule, conglomerability over the discrete
data, and sequential consistency: $P(H\mid A_{n+1})$ is the one-step update of
$P(H\mid A_n)$) holds automatically. Explicitly, the sequential update factor is
\begin{equation}\label{eq:sequential}
  \frac{P(H\mid A_{n+1})}{P(H\mid A_n)}
  =\frac{P(A_n)}{P(A_{n+1})}
  =\frac{p+(1-p)m_n}{p+(1-p)m_{n+1}}\ \ge 1,
\end{equation}
which is exactly Bayes' rule applied to the single new observation
$X_{n+1}=1$ (the inequality holds because $m_{n+1}\le m_n$), confirming that the
increasing posterior of Theorem~\ref{thm:main}(ii) is produced by lawful
conditioning and not by any extra-Bayesian device.

\subsection{Invariance under reparametrization}
The previous treatments of this problem are vulnerable to the objection that one
simply \emph{declares} the point $\theta=1$ special and only admits relabellings
that fix it. We avoid this by first proving that $\theta=1$ is distinguished
\emph{intrinsically}, by a property of the statistical family
$\{P_\theta\}_{\theta\in[0,1]}$ alone --- with no reference to the coordinate
$\theta$ --- and then \emph{deriving} (rather than imposing) the fact that every
admissible relabelling fixes it.

\begin{lemma}[The success-endpoint is intrinsic]\label{lem:intrinsic}
Recall $A_\infty=\{X_i=1\ \forall i\}$, a tail event of the path space
$\{0,1\}^\N$ defined without reference to any parametrization. Then
\[
  P_\theta(A_\infty)=\begin{cases}1,&\theta=1,\\[1mm]0,&\theta\in[0,1),\end{cases}
  \qquad\text{so}\qquad
  \{\theta\in[0,1]:P_\theta(A_\infty)=1\}=\{1\}.
\]
Hence $\theta=1$ is the \emph{unique} parameter value whose law assigns
probability one to the tail event ``all observations are successes.''
\end{lemma}

\begin{proof}
$A_\infty=\bigcap_n A_n$ and $A_{n+1}\subseteq A_n$, so by downward continuity
$P_\theta(A_\infty)=\lim_n P_\theta(A_n)=\lim_n\theta^n$. For $\theta=1$ this
limit is $1$; for $\theta\in[0,1)$ it is $0$. The displayed set equality is
immediate.
\end{proof}

\paragraph{Reparametrizations as model isomorphisms.}
A reparametrization is a relabelling of the parameter by a Borel isomorphism
$\phi:[0,1]\to[0,1]$; we write $\eta=\phi(\theta)$ and $\psi:=\phi^{-1}$, so
$\theta=\psi(\eta)$. The defining requirement of a reparametrization is that it
leave the \emph{statistical model itself} unchanged: it merely renames the index,
so the law indexed by $\eta$ in the new chart must be the \emph{same} law that was
indexed by $\psi(\eta)$ in the old chart,
\begin{equation}\label{eq:modeliso}
  \tilde P_\eta:=P_{\psi(\eta)}\qquad(\eta\in[0,1]),
\end{equation}
and the prior is transported by push-forward, $\tilde\mu:=\phi_\#\mu$. In
particular the likelihood of a success run becomes
$\tilde P_\eta(A_n)=P_{\psi(\eta)}(A_n)=\psi(\eta)^n$. Crucially we do
\emph{not} postulate anything about $\phi(1)$; the next lemma shows the relevant
fact is \emph{forced}.

\begin{lemma}[The universal hypothesis transports to the intrinsic
hypothesis]\label{lem:fix1}
Let $\phi:[0,1]\to[0,1]$ be any Borel isomorphism and $\tilde P_\eta$ defined
by~\eqref{eq:modeliso}. Then the push-forward $\tilde H:=\phi(H)=\{\eta=\phi(1)\}$
of the universal hypothesis is exactly the intrinsically distinguished hypothesis
in the new chart:
\[
  \tilde H=\{\eta:\tilde P_\eta(A_\infty)=1\}.
\]
That is, ``the universal hypothesis'' refers, in every chart, to the unique law
charging the tail event $A_\infty$; no point is privileged by fiat. If moreover
the new chart, like the old, indexes each law by its per-trial success
probability $\Pr(X_1=1)$ (the standard convention), then $\phi(1)=1$.
\end{lemma}

\begin{proof}
$A_\infty\subseteq\{0,1\}^\N$ is one fixed set, so ``assigns $A_\infty$
probability one'' is a property of a law. By Lemma~\ref{lem:intrinsic} the unique
such law in the family is $P_1=Q_\star$. Its index in the new chart is
\[
  \{\eta:\tilde P_\eta(A_\infty)=1\}
  =\{\eta:P_{\psi(\eta)}(A_\infty)=1\}
  =\{\eta:\psi(\eta)=1\}=\{\phi(1)\}=\phi(H),
\]
using~\eqref{eq:modeliso} and Lemma~\ref{lem:intrinsic}; this is the first claim.
Hence the universal hypothesis is identified across charts not because we fixed a
point but because it is the chart-free statement ``the all-successes tail event
has probability one.'' Finally, if both charts index by the intrinsic functional
$\Pr(X_1=1)$, then $Q_\star$, for which $\Pr(X_1=1)=1$, has index $1$ in each, so
$\phi(1)=1$.
\end{proof}

Thus the universal hypothesis is \emph{not} fixed by fiat: in every chart it is
the chart-free statement ``the all-successes tail event has probability one,''
which Lemma~\ref{lem:fix1} identifies with $\phi(H)$. With this identification the
invariance computation goes through for \emph{every} Borel isomorphism, with no
admissibility restriction whatsoever.

\begin{proposition}[Reparametrization invariance]\label{prop:invariance}
Let $\phi:[0,1]\to[0,1]$ be \emph{any} Borel isomorphism, $\psi=\phi^{-1}$, and
$\tilde P$ the model in the new chart (prior $\tilde\mu=\phi_\#\mu$, likelihood
$\tilde P_\eta(A_n)=\psi(\eta)^n$). Write $\tilde H:=\phi(H)=\{\eta=\phi(1)\}$ for
the transported universal hypothesis, which by Lemma~\ref{lem:fix1} equals the
intrinsic hypothesis $\{\eta:\tilde P_\eta(A_\infty)=1\}$ in the new chart.
\begin{enumerate}
\item[\textup{(a)}] \textup{(Class is closed.)} If $\mu=p\delta_1+(1-p)\nu$ is a
spike-and-slab prior then $\tilde\mu=p\,\delta_{\phi(1)}+(1-p)\,\tilde\nu$ with
$\tilde\nu=\phi_\#\nu$ and $\tilde\nu(\{\phi(1)\})=0$; so $\tilde\mu$ has an atom
of the \emph{same} weight $p$ at the new index $\phi(1)$ of the universal
hypothesis, and is again of spike-and-slab type relative to $\tilde H$. Under the
standard success-probability convention $\phi(1)=1$ and this reads
$\tilde\mu=p\delta_1+(1-p)\tilde\nu$.
\item[\textup{(b)}] \textup{(Posterior is invariant.)} For every $n$,
$\ \tilde P(\tilde H\mid A_n)=P(H\mid A_n).$
\end{enumerate}
\end{proposition}

\begin{proof}
(a) Push-forward of a sum is the sum of push-forwards, and
$\phi_\#\delta_1=\delta_{\phi(1)}$. As $\phi$ is a bijection,
$\tilde\nu(\{\phi(1)\})=(\phi_\#\nu)(\{\phi(1)\})=\nu(\psi(\{\phi(1)\}))
=\nu(\{1\})=0$. The final sentence uses Lemma~\ref{lem:fix1}.

(b) By the image-measure formula $\int g\,d(\phi_\#\mu)=\int(g\circ\phi)\,d\mu$
and $\psi\circ\phi=\mathrm{id}$,
\[
  \tilde P(A_n)=\int_{[0,1]}\psi(\eta)^n\,d(\phi_\#\mu)(\eta)
  =\int_{[0,1]}\psi(\phi(\theta))^n\,d\mu(\theta)
  =\int_{[0,1]}\theta^n\,d\mu(\theta)=P(A_n).
\]
Since $\tilde H=\{\eta=\phi(1)\}$ and $\psi(\phi(1))=1$,
\[
  \tilde P(\tilde H\cap A_n)=\int_{\{\phi(1)\}}\psi(\eta)^n\,d\tilde\mu(\eta)
  =\psi(\phi(1))^n\,\tilde\mu(\{\phi(1)\})=1^n\cdot p=p=P(H\cap A_n),
\]
where $\tilde\mu(\{\phi(1)\})=(\phi_\#\mu)(\{\phi(1)\})=\mu(\{1\})=p$. Dividing,
$\tilde P(\tilde H\mid A_n)=p/P(A_n)=P(H\mid A_n)$.
\end{proof}

\begin{remark}[Division of labour: what invariance settles, and what it does
not]\label{rem:invariance}
It is important to state precisely which desideratum invariance discharges, since
invariance \emph{by itself} neither produces nor justifies the atom.
\begin{enumerate}
\item[(I)] \emph{Invariance does not force an atom.} The Jeffreys / Fisher\nobreakdash-information
prior $d\theta/\sqrt{\theta(1-\theta)}$ is exactly reparametrization-invariant
(its defining property), yet it is absolutely continuous and has $\mu(\{1\})=0$;
by Proposition~\ref{prop:obstruction} it cannot confirm $H$. So one cannot derive
the spike from invariance. What \emph{forces} the atom is \emph{confirmability}:
Theorem~\ref{thm:converse} shows that an exchangeable model confirms $H$
\emph{iff} its de~Finetti measure has an atom at $\theta=1$. Confirmability is the
operative principle; invariance is a consistency requirement the solution must
\emph{also} meet, not the source of the atom.
\item[(II)] \emph{What invariance does settle} is that the atom is a
\emph{coordinate-free} object, so that meeting requirement~(iii) does not
re-introduce, in disguise, the chart-dependence that caused the obstruction. By
Lemma~\ref{lem:fix1} the relevant point $\theta=1$ is intrinsically fixed by the
tail event $A_\infty$; by Proposition~\ref{prop:invariance} the atom's weight
$p=\mu(\{1\})$ and the whole posterior trajectory $\{P(H\mid A_n)\}_n$ are
invariants. By contrast ``$\mu\ll\mathrm{Leb}$'' is \emph{not} coordinate-free:
no Borel isomorphism can turn a Lebesgue-null point into one of positive Lebesgue
measure, so the obstruction's hypothesis is itself the chart-dependent artifact.
The combined statement is: \emph{confirmability selects the atom; invariance
certifies that doing so is legitimate, intrinsic, and parametrization-free.}
\item[(III)] \emph{The weight $p$ is still free at this stage.}
Proposition~\ref{prop:invariance} fixes $p$ only up to the choice of slab; any
$p\in(0,1)$ yields an invariant, confirming model. A canonical, principled value
of $p$ requires an additional selection principle, supplied in
\S\ref{sec:mdl}--\S\ref{sec:solomonoff} by a description-length / Solomonoff
argument.
\end{enumerate}
\end{remark}


\subsection{Predictive behaviour in the non-universal regime}
We verify that away from the pure-success path the model behaves like an
ordinary, well-calibrated continuous Bayesian learner, and in particular
recovers the Bayes--Laplace rule of succession.

\begin{proposition}[Reduction to standard predictive behaviour]\label{prop:predictive}
Use the spike-and-slab prior~\eqref{eq:prior} with a slab $\nu$ that charges
$(0,1)$.
\begin{enumerate}
\item[\textup{(a)}] \emph{(Falsification.)} If at any stage some $X_i=0$ is
observed, then on that event the posterior of $H$ is $0$ and stays $0$ forever;
formally $P(H\mid X_1=x_1,\dots,X_n=x_n)=0$ whenever $s_n<n$.
\item[\textup{(b)}] \emph{(Continuous learning off the spike.)} Conditional on
any data string containing a $0$, the posterior over $\theta$ is the ordinary
Bayesian posterior of the slab model $\nu$; the spike contributes nothing.
\item[\textup{(c)}] \emph{(Rule of succession; recovery of Bayes--Laplace.)}
With the uniform slab $\nu=\mathrm{Unif}[0,1]$ one has $m_n=\frac1{n+1}$ and
\[
  P(X_{n+1}=1\mid A_n)=\frac{p+(1-p)\frac{1}{n+2}}{\,p+(1-p)\frac{1}{n+1}\,}.
\]
As $p\downarrow0$ this equals $\tfrac{n+1}{n+2}$, Laplace's rule of succession,
for every $n$; for fixed $p\in(0,1)$ it lies strictly above $\tfrac{n+1}{n+2}$
and increases to $1$, the (mild) optimism induced by the spike.
\end{enumerate}
\end{proposition}

\begin{proof}
(a) If $s_n<n$ then $x_j=0$ for some $j\le n$, and the spike's likelihood is
$P_1(X_1=x_1,\dots,X_n=x_n)=1^{s_n}(1-1)^{n-s_n}=0$ because $n-s_n\ge1$. Thus
$P(H\cap\{X_1=x_1,\dots,X_n=x_n\})=p\cdot0=0$, while the slab gives the data
positive probability because $\nu$ charges $(0,1)$, where
$\theta^{s_n}(1-\theta)^{n-s_n}>0$; so the posterior of $H$ is
$0/(\text{positive})=0$. As $H$ has posterior $0$, all further conditioning keeps
it at $0$: each subsequent $A_{n+k}$-type conditioning event still gives the
spike likelihood $0$, so the spike never re-enters the posterior. (We do not claim
the general impossibility of reviving null events; the statement is specific to
this likelihood structure.)

(b) On data with $s_n<n$ the spike has zero likelihood by (a), so in the
posterior decomposition
$d\mu_{\text{post}}(\theta)\propto\theta^{s_n}(1-\theta)^{n-s_n}\,d\mu(\theta)$
the atom at $1$ receives weight $0$ and the surviving mass is
$\propto\theta^{s_n}(1-\theta)^{n-s_n}\,d\nu(\theta)$, which is precisely the
posterior obtained from the slab prior $\nu$ alone. Hence all predictions
coincide with those of the continuous model $\nu$.

(c) For $\nu=\mathrm{Unif}[0,1]$, $m_n=\int_0^1\theta^n\,d\theta=\frac1{n+1}$.
By the chain rule and Lemma~\ref{lem:marginal},
\[
  P(X_{n+1}=1\mid A_n)=\frac{P(A_{n+1})}{P(A_n)}
  =\frac{p+(1-p)\frac{1}{n+2}}{p+(1-p)\frac{1}{n+1}} .
\]
Setting $p=0$ gives $\frac{1/(n+2)}{1/(n+1)}=\frac{n+1}{n+2}$, Laplace's rule.
For $p\in(0,1)$: writing the predictive value as $f(p)$, one has $f(0)=\frac{n+1}{n+2}$
and $f(1)=1$; since $\frac1{n+1}>\frac1{n+2}>0$, cross-multiplication shows
$p+(1-p)\frac1{n+2}$ and $p+(1-p)\frac1{n+1}$ are both increasing in $p$ but the
ratio $f(p)$ is strictly increasing in $p$ on $[0,1]$ (its derivative has the sign
of $\frac1{n+1}-\frac1{n+2}>0$), so $f(p)>\frac{n+1}{n+2}$ for $p\in(0,1)$; and
as $n\to\infty$ both numerator and denominator tend to $p$, so $f(p)\to1$.
\end{proof}

Parts (a)--(b) deliver the philosophically essential feature: the spike is a
\emph{falsifiable} hypothesis in the Popperian sense---a single counterexample
($X_i=0$) annihilates its posterior---yet it accrues confirmation under a
growing run of confirming instances. This is exactly the asymmetry one wants for
``all ravens are black'': one non-black raven refutes it; many black ravens
progressively confirm it.

\section{Characterization: a converse}\label{sec:converse}

We now show that the spike is not merely \emph{sufficient} but \emph{necessary}:
within the entire class of exchangeable Bernoulli models, confirmability of $H$
in the sense of Theorem~\ref{thm:main}(ii) is equivalent to the presence of an
atom at $\theta=1$. This is what makes the admissible class a genuine
\emph{characterization}.

\begin{theorem}[Confirmability $\iff$ atom at the success-endpoint]\label{thm:converse}
Let $\mu$ be any prior on $[0,1]$, with associated exchangeable model $P$
of~\eqref{eq:joint}, and suppose $\mu\ne\delta_0$ so that $P(A_n)>0$ for all $n$.
Set $p:=\mu(\{1\})$. Then for all $n$
\begin{equation}\label{eq:general}
  P(H\mid A_n)=\frac{p}{\int_{[0,1]}\theta^n\,d\mu(\theta)} ,
\end{equation}
and consequently:
\begin{enumerate}
\item[\textup{(a)}] $P(H\mid A_n)>0$ for some (equivalently every) $n$ $\iff$
$p>0$;
\item[\textup{(b)}] $\displaystyle\lim_{n\to\infty}P(H\mid A_n)=1 \iff p>0$.
\end{enumerate}
In particular, an exchangeable Bernoulli model confirms the universal hypothesis
$H$ (in the sense of property (ii): posterior strictly positive and converging to
$1$ along a success run) \emph{if and only if} its de~Finetti mixing measure
$\mu$ places strictly positive mass on $\theta=1$. Every such $\mu$ with $p<1$
is a spike-and-slab prior (Definition~\ref{def:spikeslab}) with slab
$\nu=(\mu-p\delta_1)/(1-p)$, and $\mu=\delta_1$ when $p=1$.
\end{theorem}

\begin{proof}
By~\eqref{eq:likelihood}, $P(A_n)=\int_{[0,1]}\theta^n\,d\mu(\theta)$, and
$P(H\cap A_n)=\int_{\{1\}}\theta^n\,d\mu=p$, which gives~\eqref{eq:general}; the
denominator is positive since $\mu\ne\delta_0$ implies $\mu$ charges $(0,1]$,
where $\theta^n>0$.

(a) If $p=0$ then the numerator of~\eqref{eq:general} is $0$, so
$P(H\mid A_n)=0$ for every $n$. If $p>0$ the numerator is positive and the
denominator finite and positive, so $P(H\mid A_n)>0$ for every $n$. (The
``some/every'' equivalence holds because the numerator does not depend on $n$.)

(b) If $p=0$ then $P(H\mid A_n)\equiv0\not\to1$. Suppose $p>0$. If $p=1$ then
$\mu=\delta_1$, $P(A_n)=1$ and $P(H\mid A_n)=1$ for all $n$. If $0<p<1$, write
$\mu=p\,\delta_1+(1-p)\nu$ with $\nu=(\mu-p\delta_1)/(1-p)$, a probability
measure with $\nu(\{1\})=(\mu(\{1\})-p)/(1-p)=0$. Then
$\int_{[0,1]}\theta^n\,d\mu=p+(1-p)m_n$ with $m_n=\int_{[0,1)}\theta^n\,d\nu\to0$
by dominated convergence exactly as in Lemma~\ref{lem:marginal}. Hence
$P(H\mid A_n)=p/(p+(1-p)m_n)\to p/p=1$. The displayed iff statements combine (a)
and (b).

The identification of $\mu$ with the unique mixing measure of an exchangeable law
is the classical de~Finetti representation theorem for $\{0,1\}$-valued
exchangeable sequences; uniqueness of $\mu$ makes ``$\mu(\{1\})>0$'' a
well-defined, model-intrinsic criterion, independent of any parametrization.
\end{proof}

Theorem~\ref{thm:converse} is the precise sense in which the spike-and-slab
class is \emph{the} answer: no exchangeable Bayesian model outside this class can
confirm $H$, and every model inside it does, with the convergence rate governed
entirely by the slab through $m_n$ (faster mass concentration of $\nu$ near $1$
gives slower confirmation, and vice versa). The continuous Bayes--Laplace prior
fails precisely and only because it is the degenerate case $p=0$.

\section{Principled selection of the spike weight}\label{sec:mdl}

Remark~\ref{rem:invariance}(III) noted that confirmability and invariance fix the
\emph{form} of the prior (an atom at $\theta=1$) but leave its \emph{weight}
$p\in(0,1)$ free. We now supply a selection principle that fixes $p$
non\nobreakdash-arbitrarily, by a description-length (minimum description length,
MDL) argument; \S\ref{sec:solomonoff} shows the same value emerges from the
Solomonoff prior, so the two threads agree.

\paragraph{The selection principle.}
A prior over hypotheses should weight a hypothesis $h$ by $2^{-L(h)}$, where
$L(h)$ is the length, in bits, of a prefix-free code describing $h$ (Rissanen's
MDL / the universal prior of Solomonoff--Levin). Apply this at the top level of a
hierarchical model with exactly two model classes:
\begin{itemize}
\item $\mathcal H_{\mathrm{univ}}$: ``the law is the single deterministic
regularity $\theta=1$'' --- one fully specified hypothesis;
\item $\mathcal H_{\mathrm{slab}}$: ``the law is i.i.d.\ Bernoulli$(\theta)$ for a
non-degenerate $\theta$,'' with a continuous prior $\nu$ over $\theta\in[0,1)$ at
the lower level.
\end{itemize}
Let $\ell_{\mathrm{univ}}=L(\mathcal H_{\mathrm{univ}})$ and
$\ell_{\mathrm{slab}}=L(\mathcal H_{\mathrm{slab}})$ be the code lengths of the
two top-level class labels. The induced top-level weights, after normalization,
are
\begin{equation}\label{eq:pMDL}
  p=\frac{2^{-\ell_{\mathrm{univ}}}}{2^{-\ell_{\mathrm{univ}}}+2^{-\ell_{\mathrm{slab}}}}
   =\frac{1}{1+2^{-(\ell_{\mathrm{slab}}-\ell_{\mathrm{univ}})}} .
\end{equation}
This is a determinate value of $p$ once a code is fixed, and it is
\emph{principled} in the precise sense that it is forced by the same universal
coding rule that governs the slab. Two features matter.
\begin{enumerate}
\item[(a)] \emph{Simplicity bias (Occam / Jeffreys' simplicity postulate).} The
universal hypothesis is the unique deterministic, parameter-free regularity
``always $1$,'' whereas the slab class additionally requires naming a continuous
parameter. Under any reasonable prefix code the deterministic regularity is at
least as short to describe, $\ell_{\mathrm{univ}}\le\ell_{\mathrm{slab}}$, so
$p\ge\tfrac12$: the framework is \emph{a priori} biased toward the simpler,
exception-free generalization. This is exactly Jeffreys' simplicity postulate
(simpler laws get higher prior), here \emph{derived} from coding rather than
posited.
\item[(b)] \emph{Machine-invariance.} Replacing the reference code changes each
$\ell$ by at most an additive constant (the Invariance Theorem,
Lemma~\ref{lem:fix1}'s algorithmic analogue used in \S\ref{sec:solomonoff}), so
$p$ changes by at most a bounded multiplicative factor and, crucially,
\emph{remains in $(0,1)$}; thus $p$ is pinned down up to the irreducible $O(1)$
ambiguity inherent in all universal coding, never to $0$. The qualitative
conclusions of Theorem~\ref{thm:main} (positive prior, monotone posterior,
limit $1$) are completely insensitive to this residual ambiguity.
\end{enumerate}

\paragraph{Why this answers the ``free parameter'' objection.}
The objection to a hand-chosen $p$ is that nothing in the problem singles out a
value. The MDL principle does single one out --- \eqref{eq:pMDL} --- and does so
by the \emph{same} rule that one must already adopt to make algorithmic induction
well defined. The only residual freedom is the universal machine, which affects
$p$ by a bounded factor and cannot make $p$ vanish. There is therefore no
\emph{additional} arbitrary choice beyond the unavoidable, and universally
acknowledged, choice of reference machine. In the limit of an infinitely
expressive slab the construction passes continuously to the fully discrete
Solomonoff prior of \S\ref{sec:solomonoff}, where $p$ becomes literally
$2^{-K(Q_\star)}/Z$.

\section{The Solomonoff / algorithmic-probability realization}\label{sec:solomonoff}

We now show that the same resolution arises canonically, without an \emph{ad hoc}
choice of spike weight, when one replaces the parametric prior by a universal
(Solomonoff) prior on hypotheses; this answers the algorithmic-probability part
of the problem and explains \emph{why} algorithmic induction is immune to the
obstruction.

\subsection{Construction}
Work directly on a countable space of \emph{hypotheses} (computable measures on
$\{0,1\}^\N$) rather than on $[0,1]$. The universal hypothesis is the
deterministic ``all successes'' law
\begin{equation}\label{eq:Qstar}
  Q_\star:=\delta_{(1,1,1,\dots)},\qquad Q_\star(A_n)=1\ \forall n,\quad
  Q_\star(\{X_i=0\})=0 .
\end{equation}
Fix a prefix-free universal description, let $K(\cdot)$ be prefix Kolmogorov
complexity, and for a countable class $\mathcal M=\{Q_1,Q_2,\dots\}$ of
pairwise-distinct computable measures with $Q_\star\in\mathcal M$ define the
Solomonoff-style mixture
\begin{equation}\label{eq:Mprior}
  M(\cdot)=\sum_{k\ge1} w_k\,Q_k(\cdot),\qquad
  w_k:=\frac{2^{-K(Q_k)}}{Z},\quad Z:=\sum_{j\ge1}2^{-K(Q_j)}\in(0,1],
\end{equation}
so each $w_k>0$, $\sum_k w_k=1$, and $w_\star>0$ is the weight of $Q_\star$. The
posterior of the hypothesis $Q_\star$ is ordinary discrete Bayes,
\begin{equation}\label{eq:solpost}
  M(Q_\star\mid A_n)=\frac{w_\star\,Q_\star(A_n)}{\sum_k w_k\,Q_k(A_n)}
  =\frac{w_\star}{\sum_k w_k\,Q_k(A_n)} .
\end{equation}

\paragraph{The role of the tail event $A_\infty$.}
The behaviour of~\eqref{eq:solpost} as $n\to\infty$ is governed entirely by which
hypotheses charge the tail event $A_\infty=\bigcap_n A_n$. For each $k$,
$Q_k(A_n)\downarrow Q_k(A_\infty)$, so the denominator tends to
$\sum_k w_k Q_k(A_\infty)$ and
\begin{equation}\label{eq:sollimit}
  \lim_{n\to\infty}M(Q_\star\mid A_n)=\frac{w_\star}{\sum_k w_k\,Q_k(A_\infty)} .
\end{equation}
This raises a genuine subtlety, which we address head-on: the \emph{full}
universal class --- all computable measures --- contains laws other than
$Q_\star$ that charge $A_\infty$, e.g.\
\[
  Q^{\dagger}:=\tfrac12\,\delta_{(1,1,1,\dots)}+\tfrac12\,\delta_{(0,0,0,\dots)},
  \qquad Q^{\dagger}(A_\infty)=\tfrac12>0,
\]
which is computable and distinct from $Q_\star$. For the \emph{unrestricted}
mixture, \eqref{eq:sollimit} is then \emph{strictly less than $1$}: confirmation
is positive and monotone increasing but bounded away from certainty. So the naive
claim ``the Solomonoff posterior of the universal hypothesis tends to $1$'' is
\emph{false} for the full class, and we do not assert it.

\paragraph{The natural admissible class.}
The cure is to identify the \emph{right} hypothesis class, not by fiat but by a
structural principle: the universal hypothesis is a statement about the
\emph{generating mechanism per trial}, namely ``every trial succeeds with
certainty.'' The mechanisms of this kind are exactly the i.i.d.\ laws, i.e.\ the
exchangeable extreme points $P_\theta$ of de~Finetti's theorem. We therefore take
\begin{equation}\label{eq:Mbern}
  \mathcal M_{\mathrm{Bern}}
  :=\{\,P_\theta:\theta\in[0,1]\ \text{computable}\,\},
\end{equation}
the computable i.i.d.\ Bernoulli laws, with $Q_\star=P_1\in\mathcal M_{\mathrm{Bern}}$.
This class is \emph{natural} for three reasons, none ad hoc:
\begin{enumerate}
\item[(N1)] It is the algorithmic restriction of \emph{exactly the model in the
problem statement} (i.i.d.\ Bernoulli trials); the parametric and algorithmic
treatments are then the same object viewed two ways, the slab being replaced by a
computable dense set of $\theta$'s.
\item[(N2)] By Lemma~\ref{lem:intrinsic}, among i.i.d.\ laws \emph{only}
$\theta=1$ charges $A_\infty$: $P_\theta(A_\infty)=0$ for every $\theta<1$. The
pathological $Q^{\dagger}$ above is \emph{not} i.i.d.\ (it is a mixture of two
point masses, hence exchangeable-but-not-extreme), so it is excluded
\emph{automatically} by the per-trial reading of the universal hypothesis, not by
hand-removal.
\item[(N3)] It is closed under the only invariance at issue
(Lemma~\ref{lem:fix1}): reparametrization permutes the labels $\theta$ and fixes
$\theta=1$, hence permutes $\mathcal M_{\mathrm{Bern}}$ and fixes $Q_\star$.
\end{enumerate}

\begin{theorem}[Solomonoff confirmation of the universal hypothesis]\label{thm:sol}
Let $M$ be the mixture~\eqref{eq:Mprior} over a class $\mathcal M\ni Q_\star$.
\begin{enumerate}
\item[\textup{(i)}] $M$ gives the universal hypothesis strictly positive prior
weight $w_\star=2^{-K(Q_\star)}/Z>0$.
\item[\textup{(ii)}] $M(Q_\star\mid A_n)\in(0,1]$ is non-decreasing in $n$,
strictly exceeds $w_\star$ once some positively-weighted $Q_k$ has $Q_k(A_n)<1$,
and satisfies~\eqref{eq:sollimit}. Consequently:
\begin{enumerate}
\item[\textup{(ii-a)}] \emph{(General class.)} The limit equals $1$ \emph{iff}
$Q_\star$ is the only law in $\mathcal M$ that charges $A_\infty$. For the
\emph{full} computable class this fails (e.g.\ via $Q^{\dagger}$), so there the
posterior increases to a limit in $(w_\star,1)$ --- positive, monotone
confirmation, but not certainty.
\item[\textup{(ii-b)}] \emph{(Natural class~\eqref{eq:Mbern}.)} For
$\mathcal M=\mathcal M_{\mathrm{Bern}}$ one has, by (N2), $Q_k(A_\infty)=0$ for
every $Q_k\ne Q_\star$, so the denominator of~\eqref{eq:sollimit} is $w_\star$ and
\[
  \lim_{n\to\infty}M(Q_\star\mid A_n)=1.
\]
\end{enumerate}
\item[\textup{(iii)}] The scheme is coherent (discrete Bayes on $(w_k)$),
machine-independent up to a multiplicative $O(1)$ factor in the weights (hence the
limit in (ii-b) is unaffected), and off the success path it reduces to the
standard Solomonoff predictor, which is consistent.
\end{enumerate}
\end{theorem}

\begin{proof}
(i) Immediate from~\eqref{eq:Mprior}, as $Q_\star$ is computable, $K(Q_\star)<\infty$.

(ii) $Q_\star(A_n)=1$ and each $Q_k(A_n)$ is non-increasing in $n$, so
$D_n:=\sum_k w_k Q_k(A_n)$ is non-increasing and $\ge w_\star Q_\star(A_n)=w_\star>0$;
hence $M(Q_\star\mid A_n)=w_\star/D_n$ is non-decreasing in $(0,1]$ and exceeds
$w_\star$ once $D_n<1$. By monotone convergence (dominating summable sequence
$(w_k)$), $D_n\to\sum_k w_k Q_k(A_\infty)$, giving~\eqref{eq:sollimit}.
For (ii-a): since $Q_\star(A_\infty)=1$, the limit is $1$ iff
$\sum_{k:Q_k\ne Q_\star}w_k Q_k(A_\infty)=0$, i.e.\ iff no other positively
weighted law charges $A_\infty$; the witness $Q^{\dagger}$ (computable,
$Q^{\dagger}(A_\infty)=\tfrac12>0$) shows this fails for the full class, where the
limit is $w_\star/(w_\star+\sum_{k\ne\star}w_kQ_k(A_\infty))\in(w_\star,1)$.
For (ii-b): by (N2) every $P_\theta\in\mathcal M_{\mathrm{Bern}}$ with $\theta<1$
has $P_\theta(A_\infty)=\lim_n\theta^n=0$ (Lemma~\ref{lem:intrinsic}), and
$\theta=1$ gives $Q_\star$; so $\sum_k w_k Q_k(A_\infty)=w_\star$ and the limit
is $1$.

(iii) Coherence is Bayes' rule for $(w_k)$. Machine-independence is the Invariance
Theorem: for universal prefix machines $U,U'$ there is $c$ with
$|K_U(Q)-K_{U'}(Q)|\le c$ for all $Q$, so unnormalized weights differ by a factor
in $[2^{-c},2^{c}]$; this leaves $w_\star>0$ and, since in (ii-b) the limit is
$w_\star/w_\star=1$ \emph{regardless} of the common normalization, the limit is
unchanged. Off the success path some $X_i=0$ forces $Q_\star(A_n)=0$, removing
$Q_\star$; $M$ becomes the standard mixture over $\mathcal M\setminus\{Q_\star\}$,
whose predictions converge a.s.\ to the true conditional for any computable source
(Solomonoff consistency).
\end{proof}

\begin{remark}[Convergence rate and the canonical $p$]\label{rem:rate}
For $\mathcal M_{\mathrm{Bern}}$ the spike weight is fixed canonically as
$p:=w_\star=2^{-K(Q_\star)}/Z$, matching~\eqref{eq:pMDL} in the discrete limit and
removing the free parameter of \S\ref{sec:class}. The convergence rate is
controlled by $K(Q_\star)$ through $w_\star$: from~\eqref{eq:solpost},
\[
  1-M(Q_\star\mid A_n)=\frac{D_n-w_\star}{D_n}
  =\frac{\sum_{k\ne\star}w_kQ_k(A_n)}{D_n}
  \le \frac{1}{w_\star}\sum_{k\ne\star}w_kQ_k(A_n),
\]
so the shortfall from certainty decays at the rate at which the rival weighted
likelihoods $\sum_{k\ne\star}w_kQ_k(A_n)$ vanish; the leading rival is the
shortest-code $\theta<1$, contributing $\asymp 2^{-K(P_\theta)}\theta^n$, an
exponential-in-$n$ decay with prefactor set by the complexity of the simplest
alternative. A smaller $K(Q_\star)$ (simpler universal law) gives a larger
$w_\star$ and hence faster confirmation.
\end{remark}

\subsection{Why algorithmic probability resolves the obstruction
(and what it does \emph{not} give for free)}
The parametric model required a spike by hand because the continuum $[0,1]$
contains the universal hypothesis as a measure-zero point and the natural
continuous measures cannot charge it. In the algorithmic setting the hypothesis
space is \emph{countable}, so the universal prior~\eqref{eq:Mprior} is
\emph{discrete} and \emph{every} computable hypothesis, including $Q_\star$,
automatically gets positive weight $\propto2^{-K(Q)}$. Hence
\emph{positivity of the prior} (requirement (i)) and a \emph{canonical} spike
weight $p=2^{-K(Q_\star)}/Z$ are obtained for free, and the complexity weighting
realizes Occam's razor / Jeffreys' simplicity postulate in favour of the simplest
universal regularity.

What is \emph{not} free is the convergence of the posterior to certainty
(requirement (ii)). As \S\ref{sec:solomonoff} shows explicitly, over the
\emph{full} computable class the limit is $<1$, because computable laws other than
$Q_\star$ (such as $Q^{\dagger}$) also charge the tail event $A_\infty$.
Attaining limit $1$ requires choosing the class so that $Q_\star$ is the unique
charger of $A_\infty$. We argued this is achieved \emph{non-arbitrarily} by
$\mathcal M_{\mathrm{Bern}}$, the computable i.i.d.\ Bernoulli laws (N1)--(N3):
reading the universal hypothesis as a statement about the per-trial mechanism
forces the i.i.d.\ (de~Finetti-extreme) class, within which Lemma~\ref{lem:intrinsic}
guarantees that only $\theta=1$ charges $A_\infty$. This is the precise algorithmic
counterpart of the slab condition $\nu(\{1\})=0$: not an ad hoc deletion, but the
same intrinsic tail-event criterion (Lemma~\ref{lem:intrinsic}) that distinguishes
$\theta=1$ in the parametric story. Thus the algorithmic and parametric
resolutions are one resolution, with the spike weight pinned down canonically and
the convergence to certainty earned by the natural, principled restriction of the
hypothesis class.


\section{Summary: adequacy of the resolution}\label{sec:summary}

We collect the resolution and its justification.

\begin{enumerate}
\item \textbf{The obstruction is exactly absolute continuity}
(Proposition~\ref{prop:obstruction}): it depends only on $\mu(\{1\})=0$, a
chart-dependent constraint, not on any structural feature of confirmation.

\item \textbf{Admissible class.} The spike-and-slab priors
$\mu=p\delta_1+(1-p)\nu$, $p\in(0,1)$, $\nu(\{1\})=0$
(Definition~\ref{def:spikeslab}), give $P(H)=p>0$ and posterior
$P(H\mid A_n)=p/(p+(1-p)m_n)\uparrow1$
(Theorem~\ref{thm:main}); this is requirements (i)--(ii).

\item \textbf{Coherence and invariance} (requirement (iii)):
the model is an ordinary Bayesian model, so updating is coherent
(\eqref{eq:sequential}); the admissible class is closed under reparametrization
and the entire posterior trajectory is invariant
(Proposition~\ref{prop:invariance}, Remark~\ref{rem:invariance}); and the model
reduces to the Bayes--Laplace rule of succession in the non-universal regime,
with one observed counterexample falsifying $H$ outright
(Proposition~\ref{prop:predictive}).

\item \textbf{Characterization} (Theorem~\ref{thm:converse}): among
\emph{all} exchangeable Bernoulli models, confirmability of $H$ holds if and only
if the de~Finetti mixing measure has an atom at $\theta=1$. The class of item~2
is therefore not one device among many but the complete solution set.

\item \textbf{Principled spike weight} (\S\ref{sec:mdl}): the form of the prior
is forced by confirmability and invariance, but the weight $p$ is fixed
non\nobreakdash-arbitrarily by a minimum-description-length rule,
$p=2^{-\ell_{\mathrm{univ}}}/(2^{-\ell_{\mathrm{univ}}}+2^{-\ell_{\mathrm{slab}}})$
(\eqref{eq:pMDL}); the simpler, exception-free universal law receives
$p\ge\tfrac12$ (Jeffreys' simplicity postulate, here derived from coding), and the
only residual freedom (the reference machine) perturbs $p$ by a bounded factor and
cannot send it to $0$.

\item \textbf{Algorithmic realization} (Theorem~\ref{thm:sol},
Remark~\ref{rem:rate}): a Solomonoff prior over computable hypotheses is a
discrete mixture charging the ``all successes'' law $Q_\star$ with weight
$2^{-K(Q_\star)}/Z>0$, giving requirement~(i) and a \emph{canonical} spike weight
for free. Requirement~(ii) (posterior $\to1$) is \emph{not} automatic: over the
full computable class the limit is $<1$ because other computable laws (e.g.\
$Q^{\dagger}=\tfrac12\delta_{\mathbf 1}+\tfrac12\delta_{\mathbf 0}$) also charge
the tail event $A_\infty$. Limit $1$ holds precisely for the natural class
$\mathcal M_{\mathrm{Bern}}$ of computable i.i.d.\ Bernoulli laws, where
Lemma~\ref{lem:intrinsic} guarantees only $\theta=1$ charges $A_\infty$ --- the
algorithmic counterpart of $\nu(\{1\})=0$, motivated (N1)--(N3) by reading $H$ as a
per-trial mechanism, not by hand-removal. Convergence is exponential in $n$ with
rate set by $K(Q_\star)$ and the simplest rival (Remark~\ref{rem:rate}).
\end{enumerate}

\noindent\textbf{Philosophical adequacy.} The construction reconciles three
desiderata usually thought to be in tension. (a) \emph{Confirmability}: $H$ is
genuinely learnable, with monotone increasing posterior tending to certainty
under accumulating positive instances. (b) \emph{Falsifiability}: a single
counterexample drives the posterior of $H$ to $0$ irreversibly
(Proposition~\ref{prop:predictive}(a)), so $H$ remains an empirical, refutable
claim and the framework is not credulous. (c) \emph{Open-mindedness}: the slab
$\nu$ keeps positive belief on every $\theta<1$, so before decisive evidence the
model entertains the possibility that $H$ is false and predicts in a calibrated,
Laplace-like manner. The zero-prior obstruction is thus revealed as an avoidable
artifact of an over-restrictive (merely-continuous) prior class; once the
genuinely universal hypothesis is admitted as a first-class atom---whether by an
explicit spike or, canonically, through algorithmic probability---standard
Bayesian conditioning confirms it exactly as scientific practice demands.

\medskip
\noindent\textbf{Engagement with the standard objections.} The problem's final
clause asks us to justify adequacy, which requires meeting the recognized
objections, not merely restating the obstruction.
\begin{enumerate}
\item[(O1)] \emph{``Charging a Lebesgue-null point is illegitimate / denies that
$\theta$ is genuinely continuous.''} Our entitlement is not measure-theoretic but
inductive: by Lemma~\ref{lem:intrinsic} the value $\theta=1$ is \emph{not} an
arbitrary point of a continuum but the unique law assigning the
parametrization-free tail event $A_\infty$ probability one --- a qualitatively
distinguished hypothesis (``a genuine regularity holds with no exceptions''),
exactly the kind of hypothesis to which scientific inference assigns standing
prior credence. Treating $\theta$ as ``genuinely continuous'' is itself a modelling
choice (the merely-continuous prior); Theorem~\ref{thm:converse} shows that choice
is \emph{equivalent} to declaring universal generalizations unconfirmable. One is
free to make it, but then one has assumed the conclusion, not discovered it.
\item[(O2)] \emph{Countable additivity vs.\ invariance.} There is no tension here:
$\mu=p\delta_1+(1-p)\nu$ is a countably additive probability measure on $[0,1]$,
and Proposition~\ref{prop:invariance} establishes full invariance under all Borel
reparametrizations \emph{within} countable additivity. The atom is invariant
because point masses push forward to point masses and the success-endpoint is
intrinsically fixed (Lemma~\ref{lem:fix1}); no finitely-additive or improper prior
is needed. The tension that afflicts ``ignorance priors'' (e.g.\ the conflict
between a uniform prior and reparametrization) is avoided precisely because our
prior is not chosen to express ignorance but to encode the structural fact that a
no-exception regularity is a distinguished hypothesis.
\item[(O3)] \emph{The alternative confirmation notions the problem invites.} (i)
\emph{Nested events $H_\varepsilon=\{\theta\ge1-\varepsilon\}$.} These already have
positive continuous-prior mass, and under any slab charging neighbourhoods of $1$,
$P(H_\varepsilon\mid A_n)\to1$ for each fixed $\varepsilon>0$; so ``approximate''
universal claims are confirmable with no atom. But $\bigcap_\varepsilon
H_\varepsilon=H$ and $P(H\mid A_n)=\lim_{\varepsilon\downarrow0}P(H_\varepsilon\mid
A_n)$ fails to reach the strict statement $\theta=1$ without the atom --- which is
exactly why the genuine universal hypothesis, as opposed to its
$\varepsilon$-relaxations, needs the spike. Our framework subsumes both: the slab
delivers the $H_\varepsilon$ behaviour, the atom delivers the limiting strict $H$.
(ii) \emph{Bayes factors / likelihood ratios (Jeffreys, Popper--Miller).} The
likelihood ratio of $H$ against any $\theta<1$ after $A_n$ is $1/\theta^n\to\infty$,
so the \emph{evidence} favours $H$ unboundedly regardless of the prior; the atom is
what converts unbounded evidence into a posterior probability tending to $1$
rather than to a prior-dependent constant. This reconciles the Bayes-factor view
(evidence accrues) with the posterior view (belief converges), and answers the
Popper--Miller worry by locating confirmation in the likelihood ratio while
keeping a coherent posterior.
\end{enumerate}
In sum, the adequacy is not asserted but argued: the atom is \emph{forced} by
confirmability (Theorem~\ref{thm:converse}), \emph{legitimate and intrinsic} by
the tail-event characterization (Lemmas~\ref{lem:intrinsic}--\ref{lem:fix1}),
\emph{invariant} within countable additivity (Proposition~\ref{prop:invariance}),
\emph{principled in weight} by description length (\S\ref{sec:mdl}), and
\emph{convergent to certainty} on the natural class
(Theorem~\ref{thm:sol}(ii-b)), while reducing to standard predictive behaviour off
the success path (Proposition~\ref{prop:predictive}).
\end{document}
